\section{Phương pháp đo lường}\label{sec:framework}

Phần này xây dựng khung đo lường cho các cơ chế phối hợp giữa các model. Trước hết, chúng tôi
tách bài toán thành ba câu hỏi cốt lõi --- \emph{tiềm năng cải thiện}, \emph{khả năng khai
thác}, \emph{thiệt hại}: pool lời giải có chứa giá trị mới hay không, một bộ chọn khả thi có
khai thác được giá trị đó hay không, và bản thân giao thức có gây ra tổn thất trên những bài
vốn đã được giải đúng hay không. Từ ba câu hỏi này, các đại lượng định lượng được xây dựng
và sử dụng nhất quán trong các thí nghiệm tiếp theo.

\subsection{Ba câu hỏi cốt lõi}\label{sec:threeq}

Để xác định một cơ chế phối hợp có thực sự tạo ra giá trị, chúng tôi xem xét ba thành phần độc
lập:

\begin{itemize}\itemsep1pt
\item \textbf{Tiềm năng cải thiện} ($H$) --- \emph{dư địa} của pool: pool các lời giải được
      sinh ra có chứa lời giải mà model mạnh khi hoạt động đơn lẻ không tìm được hay không?
      Nếu không, dù bộ chọn hoạt động hoàn hảo, cơ chế phối hợp cũng không thể cải thiện kết
      quả.
\item \textbf{Khả năng khai thác} ($\kappa$) --- chất lượng của \emph{bộ chọn}: nếu pool thực
      sự chứa một lời giải mới, một tín hiệu khả thi mà không cần biết đáp án có thể chọn đúng
      lời giải đó hay không?
\item \textbf{Thiệt hại} ($D$): giao thức có làm hỏng những bài mà model mạnh vốn đã giải
      đúng hay không, và mức độ thiệt hại là bao nhiêu?
\end{itemize}

Do đó, một cơ chế chỉ có thể tạo ra giá trị ròng khi đồng thời có ba điều kiện: có tiềm năng
cải thiện (pool chứa lời giải mới), có khả năng khai thác nó (bộ chọn nhặt được), và thiệt hại
do giao thức gây ra nhỏ hơn phần cải thiện thu được. Ba nhánh nghiên cứu trong báo cáo tương ứng với ba thành phần này.

Khung trên chỉ đóng vai trò tổ chức các câu hỏi thực nghiệm; nó không phải là một mô hình dự
báo. Các đại lượng dùng để kiểm chứng từng thành phần được định nghĩa trong các tiểu mục tiếp
theo.

\subsection{Thí nghiệm tiếp xúc và đẳng thức phân rã}\label{sec:identity}

Để cô lập tác động của việc một model mạnh nhìn thấy lời giải do model yếu sinh ra, xét một cặp
model gồm model yếu và model mạnh, chẳng hạn Qwen-1.5B và Qwen-7B. Chúng tôi sử dụng ba ký hiệu
khác với các vai $P/S/V/A$ ở \S\ref{sec:setup} để tránh nhầm lẫn:

\begin{itemize}\itemsep1pt
\item $W$ --- model \textbf{y}ếu, tự giải bài; lời giải sinh ra được gọi là \emph{artifact}.
\item $I$ --- model mạnh, tự giải bài \textbf{đ}ộc lập và không quan sát artifact của $W$.
\item $E$ --- chính model mạnh đó, với cùng lệnh giải bài, nhưng được cung cấp thêm artifact
      của $W$ trong ngữ cảnh đầu vào.
\end{itemize}

Điểm quan trọng của thiết kế này là $I$ và $E$ sử dụng cùng một model và cùng nhiệm vụ; khác
biệt duy nhất là $E$ được nhìn thấy artifact của $W$. Vì vậy, hiệu $\acc(E)-\acc(I)$ đo trực
tiếp tác động của việc tiếp xúc với lời giải của model yếu. Đây là đại lượng có thể đo trong
triển khai thực tế mà không cần biết đáp án, ký hiệu là $\dhonest$.

Tiếp theo, để biết cơ chế tiếp xúc này có \emph{dư địa} cải thiện đến đâu, ta xét một bộ chọn lý
tưởng có quyền biết đáp án. Bộ chọn này trả lời đúng nếu $W$ đúng hoặc $E$ đúng; nói cách khác,
nó đạt mức chính xác
\[
\acc(\CEIL) = P(W \text{ đúng} \vee E \text{ đúng}).
\]
Gọi mức cải thiện tối đa này so với model mạnh chạy độc lập là $\dceil$. Khi đó:
\begin{equation}\label{eq:identity}
\dceil \;=\; \acc(\CEIL) - \acc(I) \;=\; G - L + R,
\end{equation}
trong đó ba thành phần có ý nghĩa:
\begin{align*}
G &= P(W \text{ đúng} \wedge I \text{ sai})
  &&\text{\textbf{cơ hội cải thiện} --- tính chất của \emph{cặp model}},\\
L &= P(W \text{ sai} \wedge I \text{ đúng} \wedge E \text{ sai})
  &&\text{\textbf{thiệt hại} --- do \emph{giao thức} gây ra},\\
R &= P(W \text{ sai} \wedge I \text{ sai} \wedge E \text{ đúng})
  &&\text{\textbf{phần cứu được} --- bài cả hai model ban đầu đều trượt}.
\end{align*}

Đẳng thức~\eqref{eq:identity} là một hệ quả đại số trực tiếp:
\[
P(W \vee E) - P(I) \;=\; P(W \wedge \neg I) \;-\; P(\neg W \wedge I \wedge \neg E)
\;+\; P(\neg W \wedge \neg I \wedge E).
\]
Nó cho phép phân tách mức cải thiện tối đa thành ba nguồn: cơ hội do model yếu cung cấp, thiệt
hại do việc tiếp xúc gây ra, và các trường hợp được cứu khi cả hai model ban đầu đều thất bại.

Ví dụ, trên 100 bài, giả sử có 15 bài $W$ đúng nhưng $I$ sai ($G=0{,}15$), 20 bài $W$ sai,
$I$ đúng nhưng $E$ sai ($L=0{,}20$), và 5 bài cả $W$ lẫn $I$ đều sai nhưng $E$ đúng
($R=0{,}05$). Khi đó
\[
\dceil = 0{,}15 - 0{,}20 + 0{,}05 = 0,
\]
nghĩa là ngay cả bộ chọn lý tưởng cũng không thể vượt model mạnh hoạt động đơn lẻ.

Do là một cận lý tưởng, $\dceil$ cần đáp án chuẩn và vì vậy không phải là đại lượng có thể triển
khai trực tiếp. Vai trò của nó là \textbf{sàng lọc}: nếu $\dceil<0$, ngay cả bộ chọn biết đáp án
cũng không thể cải thiện model mạnh, do đó không có lý do để tiếp tục tìm một bộ chọn khả thi
$\kappa$. Ngược lại, $\dceil\geq0$ chỉ cho biết rằng về mặt thông tin vẫn tồn tại dư địa để cải
thiện; nó không đảm bảo một bộ chọn thực tế có thể khai thác được dư địa đó.

Đẳng thức~\eqref{eq:identity} cũng được dùng như một kiểm tra tính nhất quán của sổ sách. Nó
khớp dữ liệu ở toàn bộ các cặp có lưu vết: 4/4 cặp trong khối thăm dò và 15/15 cặp trong thí
nghiệm chính, tổng cộng 19/19 cặp. Tuy nhiên, đây chỉ là kiểm tra đại số của dữ liệu, không phải
bằng chứng cho bất kỳ giả thuyết thực nghiệm nào.

Cuối cùng, cần phân biệt giới hạn của $\dceil$. Đại lượng này chỉ mô tả lớp cơ chế \emph{chọn
giữa $W$ và $E$ sau khi artifact của $W$ đã được đưa vào ngữ cảnh}. Các cơ chế quyết định
\emph{có nên cho model mạnh tiếp xúc với artifact hay không} là một lớp bài toán khác và có cận
riêng, được phân tích ở \S\ref{sec:kappa}.

\subsection{Hai lớp giao thức và biến ``chênh lệch năng lực''}\label{sec:protocols}

Từ phân rã trên, một khác biệt quan trọng giữa các giao thức là liệu chúng có được phép tạo ra
một lời giải mới hay chỉ được lựa chọn trong số các lời giải đã có. Chúng tôi phân giao thức
thành hai lớp dựa trên tiêu chí vận hành này: \emph{đầu ra cuối có bắt buộc phải là một phần tử
của pool đã sinh hay không}.

\begin{itemize}\itemsep1pt
\item \textbf{Tuyển chọn}: đầu ra cuối phải là một phần tử của pool, chẳng hạn bỏ phiếu đa số
      hoặc lựa chọn bằng test. Với lớp này, $L=0$ \emph{theo cấu trúc}: giao thức không thể tạo
      ra một lời giải mới sai trên một bài mà pool đã chứa lời giải đúng.
\item \textbf{Sửa chữa}: giao thức cho phép sinh thêm chuỗi mới sau khi quan sát pool, chẳng
      hạn verifier can thiệp hoặc giao thức tiếp xúc. Với lớp này, $L$ không còn bằng 0 theo
      cấu trúc và phải được đo thực nghiệm. Trong dữ liệu của chúng tôi, verifier 7B có $L$ rất
      nhỏ nhưng vẫn khác 0 (1/300 bài, \S\ref{sec:roles}), trong khi giao thức tiếp xúc có $L$
      lớn hơn đáng kể (\S\ref{sec:termL}).
\end{itemize}

Ngoài loại giao thức, một biến quan trọng khác là \textbf{chênh lệch năng lực} giữa hai model.
Chúng tôi định nghĩa \emph{chênh} là hiệu giữa độ chính xác greedy của model mạnh và model yếu
trên cùng benchmark, với độ chính xác được chuẩn hoá trên thang 0--1. Ví dụ, với cặp
1.5B$\to$7B trên MBPP, chênh lệch này là $0{,}244$.

Chênh lệch năng lực được sử dụng như biến giải thích chính trong các thí nghiệm tiếp theo. Ký
hiệu $g^\ast$ được dành cho mức chênh tại đó đường khớp của $\dceil$ đổi dấu
(\S\ref{sec:termG}). Qua đó, chúng tôi có thể khảo sát không chỉ liệu một cơ chế phối hợp có
hiệu quả hay không, mà còn \emph{điều kiện về khoảng cách năng lực giữa các model để cơ chế đó
còn có dư địa tạo giá trị}.


% =====================================================================
\section{Thiết lập thí nghiệm và quy trình}\label{sec:setup}

\todoTD{rà toàn bộ mục này}

\paragraph{Model.} Qwen2.5-Instruct 1.5B/7B/14B/32B \cite{yang2024}; Llama-3.1-8B-Instruct
\cite{grattafiori2024}; DeepSeek-Coder-6.7B-Instruct \cite{guo2024}. Sáu model này lập thành 15
cặp có hướng của thí nghiệm \S\ref{sec:termG}; các khối còn lại dùng Qwen 1.5B/7B/14B. Lượng tử
hoá đo từng model, không suy theo họ: DeepSeek-6.7B nạp nf4 hết 3,61 GB; Qwen-7B 5,21 GB;
Llama-8B và Qwen-14B không lượng tử hoá được trên bản chạy (fp16 $\sim$16/28 GB). Phần cứng:
Kaggle 2$\times$T4 (nf4/fp16) và RTX 6000 Pro (bf16). Đo được nf4 so với bf16 dịch kết quả tới
3 điểm --- cùng cỡ với ngưỡng ý nghĩa hiệu dụng ($\sim$3,3 điểm, xem Thống kê) --- nên so sánh
chéo phần cứng chỉ đọc theo \emph{dấu}, và mọi so sánh định lượng đều trong cùng (phần cứng +
độ chính xác số + bộ bài).

\paragraph{Benchmark.} GSM8K \cite{cobbe2021}; MATH-500 \cite{hendrycks2021}; MBPP
\cite{austin2021} (dải chính \texttt{task\_id} 11--510, $\sim$499 bài sau lọc; dải giữ lại
511--974, 464 bài); HumanEval \cite{chen2021}. Điểm mấu chốt cho \S\ref{sec:kappa}:
MBPP/HumanEval có cách kiểm \emph{chắc chắn} (chạy test --- lời giải qua test là bằng chứng
trực tiếp), còn MATH/GSM8K không có.

\paragraph{Vai và pipeline.} Bốn vai: $P$ (planner --- lập dàn ý), $S$ (solver --- giải), $V$
(verifier --- kiểm và được phép sửa), $A$ (aggregator --- nhận nhiều ứng viên, chọn đáp án
cuối); các tổ hợp viết liền, ví dụ \texttt{PSVA} là pipeline đủ bốn vai. Lưu ý các chữ này chỉ
\emph{vai trong pipeline}; ba nhân vật $W/I/E$ và ba số hạng $G/L/R$ ở \S\ref{sec:identity} là
bộ ký hiệu riêng của thí nghiệm tiếp xúc.

\paragraph{Bộ tổng hợp và các mốc.} Với $k$ lượt giải độc lập cho một bài:
\texttt{maj@}$k$ = lấy đáp án xuất hiện nhiều nhất (bỏ phiếu đa số);
\texttt{oracle@}$k$ = đúng nếu \emph{ít nhất một} lượt đúng (trần lý tưởng của mọi cách chọn
trong pool --- cần đáp án chuẩn, chỉ dùng làm trần);
\texttt{llm\_agg@}$k$ = một lượt LLM đọc cả $k$ ứng viên rồi chọn/tổng hợp.
\emph{Ví dụ}: 5 lượt cho đáp án $\{7, 7, 7, 12, 15\}$, đáp án chuẩn 12 --- \texttt{maj@5} chọn
7 (sai), \texttt{oracle@5} đúng (vì có một lượt ra 12). Trên code, \texttt{exec}$k$ = chạy test
với từng ứng viên, lấy ứng viên qua test (tín hiệu chắc chắn); \texttt{llm}$k$ = để LLM tự
duyệt thay vì chạy test (tín hiệu học được).

\paragraph{Đại lượng dẫn xuất.} \texttt{V\_gain} = độ chính xác pipeline có $V$ trừ pipeline
cùng cấu hình bỏ $V$; \texttt{A\_gain} tương tự cho $A$. AUC (diện tích dưới đường ROC) đo khả
năng một bộ phân loại tách lời giải đúng khỏi lời giải sai: 0,5 = đoán ngẫu nhiên, 1,0 = hoàn
hảo.

\paragraph{Quy ước đơn vị.} Độ chính xác viết ở thang 0--1; hiệu ứng viết bằng ``điểm'' = điểm
phần trăm ($+14{,}0$ điểm $= +0{,}140$); ``chênh'' cũng là hiệu accuracy, thang 0--1 (chênh
$0{,}09 \approx 9$ điểm). Chi phí có ba thang không thay thế được cho nhau: \emph{ký tự sinh}
(đo thô), \emph{token}, và \emph{FLOP} (tham số $\times$ token) --- chỉ thang FLOP so được giữa
hai cỡ model; tỷ số chi phí viết kèm dấu $\times$ (ví dụ $0{,}88\times$) để không lẫn với
accuracy.

\paragraph{Giải mã và tính tất định.} Khối luồng thông tin dùng greedy
(\texttt{do\_sample=False}): cùng cấu hình cho kết quả giống hệt từng bài (kiểm 499/499 giữa
hai tài khoản, hai ngày). Hệ quả: chạy lại cùng cấu hình không phải bằng chứng độc lập; dao
động giữa các fold đến từ \emph{khác bài}, không phải nhiễu lấy mẫu.

\paragraph{Thống kê.} Khối vai trò/chi phí dùng 5 fold rời nhau (mỗi fold 60--100 bài). Sàn
nhiễu được hiệu chuẩn bằng cách đo cùng cấu hình trên 5 fold: \texttt{V\_gain} dao động
$+1{,}0$ đến $+8{,}0$ điểm, độ lệch chuẩn giữa fold 2,65 điểm --- một phép đo đơn lẻ 100 bài
gần như không mang thông tin. Suy diễn bằng t-test trên fold ($\mathrm{SE}=s/\sqrt{k}$;
$t_{0{,}975}(4)=2{,}776$, tức ngưỡng hiệu dụng $\sim$3,3 điểm), khoảng tin cậy 95\% (KTC), và
phép thử dấu (5/5 fold cùng dấu $\Rightarrow p=0{,}031$ một phía). Khối luồng thông tin dùng
McNemar chính xác ghép cặp trên cùng bộ bài và bootstrap ghép cặp (10\,000 lần). Với 131 phép
thử toàn dự án, kỳ vọng $\sim$6,6 dương tính giả \cite{bh1995}; các ca $p \in [0{,}02;0{,}05]$
đọc dè dặt tương ứng.

\paragraph{Nhãn tin cậy và quy trình.} \emph{Mức A} = tiền đăng ký với điều kiện hợp lệ, hoặc 5
fold kèm t-test; \emph{mức B} = đo một lần hoặc phân rã hậu nghiệm --- chỉ dùng minh hoạ cơ
chế. Mọi thí nghiệm khối luồng thông tin khoá bảng diễn giải \emph{trước} khi chạy (kèm dòng
bác bỏ giả thuyết); kết quả niêm phong hash trước khi đọc; lần chạy không đạt điều kiện hợp lệ
bị đánh dấu VOID và không đọc số. Thành tích của bộ lọc này: 16/32 lần chạy VOID (50\%); sổ dự
đoán trước công khai đoán đúng 21/43 --- người đặt giả thuyết đoán sai một nửa, đó chính là lý
do phải khoá diễn giải trước khi thấy số. Trong lần rà giữa kỳ trên 131 đại lượng: bốn kết quả
chủ lực của khối fold ($+14{,}0$, $+11{,}0$, $+5{,}6$, $+4{,}4$ --- xem \S\ref{sec:results})
giữ ý nghĩa ($t=3{,}3$--$6{,}9$); kết quả $+7{,}7$ của khối McNemar cũng giữ ($p=0{,}0075$);
6 kết quả được phục hồi sau khi đối chiếu cờ hợp lệ; 1 mất ý nghĩa ($k{=}3$ fold); 2 ca treo
được phán bằng dữ liệu fold còn trên Kaggle, không cần chạy lại.

% =====================================================================
